% GENERATED FILE. Edit canonical lessons, then run npm run generate:books.
% canonical-source-hash: fnv1a64:fad20d8b1dbf74a6
% canonical-lessons: ZH-C497-zixue, ZH-C497-kewai, ZH-C497-biji, ZH-C497-zhishi, ZH-C497-shengxue, ZH-R497-review-492-497

\chapter{Learning on Your Own}
\label{ch:zh-learning-on-your-own}
\hlchaptermodality{497}

\begin{chapteropening}
\textbf{By the end of this chapter:} \emph{I can talk about teaching myself, time outside class, notes, knowledge and going on to further study.}

Complete the last lesson of chapter 497: I can talk about teaching myself, time outside class, notes, knowledge and going on to further study.
\end{chapteropening}

\section[\texorpdfstring{\zh{自学} (zìxué)}{自学 (zìxué)}]{\zh{自学} (zìxué) — to teach oneself}

\label{lesson:ZH-C497-zixue}



\begin{quote}
Before the new one: say the Mandarin for elementary (in level), then the Mandarin for to repeat a school year.
\end{quote}

\subsection*{You'll want to know: \zh{自学}}
\textbf{\zh{自学}} — \emph{zìxué} — \textquotedblleft{}to teach oneself\textquotedblright{}.

\begin{quote}
\textbf{\zh{我自学了两年汉语现在是中级}} \emph{wǒ zìxué le liǎng nián Hànyǔ xiànzài shì zhōngjí} — I taught myself Chinese for two years, and now I'm at intermediate level.
\end{quote}

\begin{cousinweb}
\textbf{\zh{自学}} (\emph{zìxué}) is built from \zh{自} 'oneself' and \zh{学} 'to learn', both inherited from Middle Chinese, \emph{dzijH} and \emph{haewk} (Guangyun, in Baxter's transcription), which glosses the first 'from'; the sense 'oneself' is not given there.
\end{cousinweb}

\begin{grammarlens}[title={What you've built}]
One more word for saying how it was.
\end{grammarlens}

\subsection*{Your turn}
\begin{itemize}
\raggedright
  \item \emph{Say it:} \emph{zìxué}
  \item \emph{Say it:} \emph{zìxué}, once more
  \item \emph{Say it:} \emph{wǒ zìxué le liǎng nián Hànyǔ xiànzài shì zhōngjí} — I taught myself Chinese for two years, and now I'm at intermediate level.
  \item \emph{Recall:} say the Mandarin for elementary (in level), then the Mandarin for to repeat a school year, then say \emph{zìxué} again
\end{itemize}

\begin{tcolorbox}[breakable,colback=teal!4,colframe=teal!35!black,title={Before you move on}]
What does \zh{自学} mean? (\textquotedblleft{}To teach oneself\textquotedblright{}.) Say it once more.
\end{tcolorbox}
\section[\texorpdfstring{\zh{课外} (kèwài)}{课外 (kèwài)}]{\zh{课外} (kèwài) — outside class; extracurricular}

\label{lesson:ZH-C497-kewai}



\begin{quote}
Before the new one: say the Mandarin for to repeat a school year, then the Mandarin for to teach oneself.
\end{quote}

\subsection*{You'll want to know: \zh{课外}}
\textbf{\zh{课外}} — \emph{kèwài} — \textquotedblleft{}outside class; extracurricular\textquotedblright{}.

\begin{quote}
\textbf{\zh{我还没有做课外作业}} \emph{wǒ hái méiyǒu zuò kèwài zuòyè} — I haven't done my homework yet.
\end{quote}

\begin{cousinweb}
\textbf{\zh{课外}} (\emph{kèwài}) is built from \zh{课} 'a lesson' and \zh{外} 'outside', both inherited from Middle Chinese, \emph{khwaH} and \emph{ngwajH} (Guangyun, in Baxter's transcription), which glosses the first 'a tax' and 'to test' and the second 'the outside'.
\end{cousinweb}

\begin{grammarlens}[title={What you've built}]
One more word for saying how it was.
\end{grammarlens}

\subsection*{Your turn}
\begin{itemize}
\raggedright
  \item \emph{Say it:} \emph{kèwài}
  \item \emph{Say it:} \emph{kèwài}, once more
  \item \emph{Say it:} \emph{wǒ hái méiyǒu zuò kèwài zuòyè} — I haven't done my homework yet.
  \item \emph{Recall:} say the Mandarin for to repeat a school year, then the Mandarin for to teach oneself, then say \emph{kèwài} again
\end{itemize}

\begin{tcolorbox}[breakable,colback=teal!4,colframe=teal!35!black,title={Before you move on}]
What does \zh{课外} mean? (\textquotedblleft{}Outside class; extracurricular\textquotedblright{}.) Say it once more.
\end{tcolorbox}
\section[bǐjì]{bǐjì — notes (taken in class)}

\label{lesson:ZH-C497-biji}



\begin{quote}
Before the new one: say the Mandarin for to teach oneself, then the Mandarin for outside class.
\end{quote}

\subsection*{You'll want to know: bǐjì}
\emph{bǐjì} — \textquotedblleft{}notes (taken in class)\textquotedblright{}.

\begin{quote}
\emph{nǐ de bǐjì xiě de hěn hǎo} — Your notes are really well written.
\end{quote}

Its characters come later; for now it is a word for the ear.

\begin{cousinweb}
\emph{bǐjì} is built from \emph{bǐ} 'a brush' and \zh{记} 'to record', both inherited from Middle Chinese, \emph{pit} and \emph{kiH} (Guangyun, in Baxter's transcription), which glosses the second 'to record'.
\end{cousinweb}

\begin{grammarlens}[title={What you've built}]
One more word for saying how it was.
\end{grammarlens}

\subsection*{Your turn}
\begin{itemize}
\raggedright
  \item \emph{Say it:} \emph{bǐjì}
  \item \emph{Say it:} \emph{bǐjì}, once more
  \item \emph{Say it:} \emph{nǐ de bǐjì xiě de hěn hǎo} — Your notes are really well written.
  \item \emph{Recall:} say the Mandarin for to teach oneself, then the Mandarin for outside class, then say \emph{bǐjì} again
\end{itemize}

\begin{tcolorbox}[breakable,colback=teal!4,colframe=teal!35!black,title={Before you move on}]
What does bǐjì mean? (\textquotedblleft{}Notes (taken in class)\textquotedblright{}.) Say it once more.
\end{tcolorbox}
\section[zhīshi]{zhīshi — knowledge}

\label{lesson:ZH-C497-zhishi}



\begin{quote}
Before the new one: say the Mandarin for outside class, then the Mandarin for notes (taken in class).
\end{quote}

\subsection*{You'll want to know: zhīshi}
\emph{zhīshi} — \textquotedblleft{}knowledge\textquotedblright{}.

\begin{quote}
\emph{zhīshi bǐ qián zhòngyào} — Knowledge matters more than money.
\end{quote}

Its characters come later; for now it is a word for the ear.

\begin{cousinweb}
\emph{zhīshi} is built from \emph{zhī} 'to know' and \emph{shí} 'to recognise', both inherited from Middle Chinese, \emph{trje} and \emph{syik} (Guangyun, in Baxter's transcription), which glosses the first 'to perceive' and the second 'to know'; in the word the second syllable is light.
\end{cousinweb}

\begin{grammarlens}[title={What you've built}]
One more word for saying how it was.
\end{grammarlens}

\subsection*{Your turn}
\begin{itemize}
\raggedright
  \item \emph{Say it:} \emph{zhīshi}
  \item \emph{Say it:} \emph{zhīshi}, once more
  \item \emph{Say it:} \emph{zhīshi bǐ qián zhòngyào} — Knowledge matters more than money.
  \item \emph{Recall:} say the Mandarin for outside class, then the Mandarin for notes (taken in class), then say \emph{zhīshi} again
\end{itemize}

\begin{tcolorbox}[breakable,colback=teal!4,colframe=teal!35!black,title={Before you move on}]
What does zhīshi mean? (\textquotedblleft{}Knowledge\textquotedblright{}.) Say it once more.
\end{tcolorbox}
\section[shēngxué]{shēngxué — to go on to a higher school; to go on to further study}

\label{lesson:ZH-C497-shengxue}



\begin{quote}
Before the new one: say the Mandarin for notes (taken in class), then the Mandarin for knowledge.
\end{quote}

\subsection*{You'll want to know: shēngxué}
\emph{shēngxué} — \textquotedblleft{}to go on to a higher school; to go on to further study\textquotedblright{}.

\begin{quote}
\emph{gāozhōng bìyè yǐhòu tā méiyǒu shēngxué} — After senior high he didn't go on to further study.
\end{quote}

Its characters come later; for now it is a word for the ear.

\begin{cousinweb}
\emph{shēngxué} is built from \emph{shēng} 'to rise' and \zh{学} 'to learn', both inherited from Middle Chinese, \emph{sying} and \emph{haewk} (Guangyun, in Baxter's transcription); the sense 'to rise' is not given for the first there.
\end{cousinweb}

\begin{grammarlens}[title={What you've built}]
One more word for saying how it was.
\end{grammarlens}

\subsection*{Your turn}
\begin{itemize}
\raggedright
  \item \emph{Say it:} \emph{shēngxué}
  \item \emph{Say it:} \emph{shēngxué}, once more
  \item \emph{Say it:} \emph{gāozhōng bìyè yǐhòu tā méiyǒu shēngxué} — After senior high he didn't go on to further study.
  \item \emph{Recall:} say the Mandarin for notes (taken in class), then the Mandarin for knowledge, then say \emph{shēngxué} again
\end{itemize}

\begin{tcolorbox}[breakable,colback=teal!4,colframe=teal!35!black,title={Before you move on}]
What does shēngxué mean? (\textquotedblleft{}To go on to a higher school; to go on to further study\textquotedblright{}.) Say it once more.
\end{tcolorbox}
\section[(dialogue)]{Looking back: from School Subjects to Learning on Your Own}

\label{lesson:ZH-R497-review-492-497}



\begin{quote}
Say the Mandarin for Chinese (as a school subject) and maths.
\end{quote}

\subsection*{Your turn}
\begin{itemize}
\raggedright
  \item \emph{Say it:} Chinese (as a school subject), maths, geography, English (the language), a major (at university)
  \item \emph{Say it:} to take extra lessons, a private tutor, to educate, to tutor, to seek advice
  \item \emph{Say it:} the college entrance exam, an exam candidate, to enter for an exam, to pass the entrance exam for, to fail an entrance exam
  \item \emph{Say it:} the middle of the term, an exam hall, full marks, to cheat in an exam, a hard problem
  \item \emph{Say it:} to learn a text by heart, advanced (in level), intermediate (in level), elementary (in level), to repeat a school year
  \item \emph{Say it:} to teach oneself, outside class, notes (taken in class), knowledge, to go on to further study
  \item \emph{Read it:} \zh{教} \zh{考} \zh{级} — each one aloud, then a word that holds it
\end{itemize}

\begin{tcolorbox}[breakable,colback=teal!4,colframe=teal!35!black,title={Before you move on}]
Chinese (as a school subject)? (\textbf{\zh{语文}}, \emph{yǔwén}.) Maths? (\emph{shùxué}.) Geography? (\textbf{\zh{地理}}, \emph{dìlǐ}.) English (the language)? (\emph{Yīngyǔ}.) A major (at university)? (\emph{zhuānyè}.) To take extra lessons? (\emph{bǔxí}.) A private tutor? (\textbf{\zh{家教}}, \emph{jiājiào}.) To educate? (\emph{jiàoyù}.) To tutor? (\emph{fǔdǎo}.) To seek advice? (\textbf{\zh{请教}}, \emph{qǐngjiào}.) The college entrance exam? (\textbf{\zh{高考}}, \emph{gāokǎo}.) An exam candidate? (\textbf{\zh{考生}}, \emph{kǎoshēng}.) To enter for an exam? (\textbf{\zh{报考}}, \emph{bàokǎo}.) To pass the entrance exam for? (\textbf{\zh{考上}}, \emph{kǎoshàng}.) To fail an entrance exam? (\emph{luòbǎng}.) The middle of the term? (\textbf{\zh{期中}}, \emph{qīzhōng}.) An exam hall? (\textbf{\zh{考场}}, \emph{kǎochǎng}.) Full marks? (\emph{mǎnfēn}.) To cheat in an exam? (\emph{zuòbì}.) A hard problem? (\textbf{\zh{难题}}, \emph{nántí}.) To learn a text by heart? (\emph{bèishū}.) Advanced (in level)? (\textbf{\zh{高级}}, \emph{gāojí}.) Intermediate (in level)? (\textbf{\zh{中级}}, \emph{zhōngjí}.) Elementary (in level)? (\emph{chūjí}.) To repeat a school year? (\emph{liújí}.) To teach oneself? (\textbf{\zh{自学}}, \emph{zìxué}.) Outside class? (\textbf{\zh{课外}}, \emph{kèwài}.) Notes (taken in class)? (\emph{bǐjì}.) Knowledge? (\emph{zhīshi}.) To go on to further study? (\emph{shēngxué}.)
\end{tcolorbox}
